% Generated draft: requires mathematical and source review.
\documentclass[11pt]{article}
\usepackage[margin=25mm]{geometry}
\usepackage{fontspec}
\setmainfont{Latin Modern Roman}
\usepackage{amsmath,amssymb,mathtools,unicode-math}
\setmathfont{Latin Modern Math}
\usepackage{enumitem,needspace,xurl,hyperref}
\setlength{\parindent}{0pt}
\setlength{\parskip}{5pt}
\setlength{\emergencystretch}{3em}
\newcommand{\N}{\mathbb{N}}
\newcommand{\Z}{\mathbb{Z}}
\newcommand{\Q}{\mathbb{Q}}
\newcommand{\R}{\mathbb{R}}
\newcommand{\C}{\mathbb{C}}
\newcommand{\F}{\mathbb{F}}
\newcommand{\A}{\mathbb{A}}
\newcommand{\PP}{\mathbb{P}}
\newcommand{\E}{\mathbb{E}}
\newcommand{\eps}{\varepsilon}
\newcommand{\dd}{\,\mathrm{d}}
\begin{document}
\section*{UnsolvedMath ID: 30004597}
\textbf{Source problem code:} catalogue entry
\subsection{Definitions and mathematical statement}
Let \(G=(V,E)\) denote a realization of a weight-dependent random connection model referred to in the supplied source. For vertices \(u,v\in V\), define their chemical distance to be their graph distance,
\[
D_G(u,v)=\inf\{k\in\Z_{\geq 0}:\text{there is a path of }k\text{ edges from }u\text{ to }v\},
\]
with \(\inf\varnothing=+\infty\). This convention defines distance for disconnected vertices; the source does not specify whether the intended asymptotic question conditions on connectivity or restricts vertices to an infinite component.

\textbf{Requested determination.} For the weight-dependent random connection models considered in the cited report, determine the asymptotic behavior of \(D_G(u,v)\) in the parameter regime described by the source as being without ultrasmallness. In particular, determine whether the reported sublogarithmic lower bound gives the sharp order of chemical distance, or whether linear or polylogarithmic growth occurs in different parameter regimes.

\textbf{Unresolved formulation qualifications.} The supplied record does not specify the model law, spatial domain, weight distribution, connection probabilities, admissible parameters, or the parameter thresholds defining the regime without ultrasmallness. It also does not specify how the two vertices are chosen, which separation or system-size variable tends to infinity, the mode of probabilistic convergence or comparison, or the formula and hypotheses of the reported lower bound. Consequently, ``sublogarithmic,'' ``linear,'' and ``polylogarithmic'' cannot be assigned precise asymptotic meanings from this record alone. The requested determination is therefore a source-qualified research question, not a fully quantified conjecture; these missing data must be recovered before a precise asymptotic assertion can be stated.
\subsection{Short English statement}
Determine how graph distance grows between two vertices in the cited weight-dependent random connection models outside the ultrasmall regime. The source asks whether a reported sublogarithmic lower bound is sharp or whether linear or polylogarithmic growth occurs in some parameter regimes. The model, limiting variable, vertex selection, and exact lower bound are missing, so a fully quantified statement requires further source review.
\subsection{Sources}
\begin{itemize}
\item Supplied local UnsolvedMath registry record, permanent ID \(30004597\), source problem code \texttt{catalogue entry}, titled \emph{Chemical Distances in Weighted Random Connection Models}, in the Oberwolfach Reports Open Problems collection. The record gives proposed year 2021 and identifies no individual proposer.
\item Supplied report citation: \emph{Spatial Networks and Percolation (2022)}, Oberwolfach Reports, pp.\ 149--224, DOI \texttt{10.4171/owr/2021/4}. \url{https://doi.org/10.4171/owr/2021/4}. The supplied literature list labels this reference ``Oberwolfach Report 4/2021.'' The record describes Oberwolfach Reports as documenting workshops and their open-problem discussions.
\item Supplied literature reference: Grauer--Gracar--Moerters, \emph{Communications in Mathematical Physics}. \url{https://doi.org/10.1007/s00220-022-04445-3}. No article title or theorem statement is supplied.
\item Literature assessment dated exactly \texttt{2026-08-22}: the supplied curation labels the problem open, states that the chemical-distance limit theorem is proved in the ultrasmall regime, and describes the non-ultrasmall case and the conjectured sharp sublogarithmic lower bound as unresolved. This is supplied dated curation, not an independently verified assessment. The background expressly requests independent verification of time-sensitive status and cited claims.
\item The registry supplies the statement status \texttt{corrected\_verified} and says that the repaired standalone statement was checked against the official Oberwolfach report. This records the supplied verification claim; no independent source consultation or mathematical verification is asserted here.
\end{itemize}
\end{document}
